\begin{lemma}
Assume the hierarchy \(\noderef{ParameterHierarchy}\) for
\((\varepsilon,\varepsilon_1,\varepsilon_2,n_0)\).  Under the two-bite
probability law with
\[
M(n)=\noderef{TwoBiteNaturalReducedVertexCount}(n),\qquad
p(n)=\noderef{TwoBiteEdgeProbability}\!\left(\frac12,n\right),
\qquad
K(n)=\noderef{TwoBiteNaturalIndependenceScale}(1+\varepsilon,n),
\]
with high probability, for every \(K(n)\)-set \(I\),
\[
\left|\bigcup_{v\in M_I}\binom{X_v(I)}2\right|\le \varepsilon_1 K(n)^2.
\]
\end{lemma}
\begin{proof}
For each \(n\), write
\[
M=M(n),\qquad K=K(n),\qquad
p=p(n)=\frac12\sqrt{\frac{\log n}{n}},
\]
and put
\[
t_2=\noderef{TwoBiteLargeCutoff}(\varepsilon,n)=n^{1/4+\varepsilon},
\qquad
t_3=\noderef{TwoBiteSmallCutoff}(\varepsilon,n)=n^{2\varepsilon}.
\]
The hypotheses of the Lean statement identify the abstract functions and edge
parameter with these quantities at each \(n\).  Namely,
\[
m(n)=M,\qquad k(n)=K,\qquad
\noderef{TwoBiteEdgeProbability}(\beta,n)
=\noderef{TwoBiteEdgeProbability}\!\left(\frac12,n\right)=p,
\]
because \(\beta=\frac12\).  Therefore the probability appearing in the Lean
conclusion is exactly
\[
\noderef{TwoBiteEventProbability}(n,M,p)
  \left(\omega\mapsto
  \noderef{TwoBiteMediumClosedPairsEvent}(K,\varepsilon,\varepsilon_1,\omega)
  \right).
\]
Thus proving the event below under the rounded \(M,K\) and the
\(\beta=\frac12\) edge law proves the displayed
\(\noderef{WithHighProbability}\) statement.

Let \(\mathcal M_n\) be the event
\(\noderef{TwoBiteMediumClosedPairsEvent}(K,\varepsilon,\varepsilon_1)\).
Unfolding \(\noderef{TwoBiteMediumClosedPairsEvent}\),
\(\noderef{ClosedPairsBound}\), and
\(\noderef{TwoBiteClosedPairsInClass}\), this event says that every
\(I\subseteq\mathrm{Fin}(n)\) with \(|I|=K\) satisfies
\[
\left|\bigcup_{x\in M_I}\binom{X_x(I)}2\right|
\le \varepsilon_1K^2, \tag{1}
\]
where \(M_I=\noderef{TwoBiteMediumClass}(\omega,\varepsilon,I)\).
Unfolding \(\noderef{TwoBiteMediumClass}\), each \(x\in M_I\) satisfies
\[
t_3<|X_x(I)|\le t_2. \tag{2}
\]

We prove that \(\Pr(\mathcal M_n^c)\to0\).  Since
\(\noderef{WithHighProbability}\) is the assertion that these event
probabilities tend to \(1\), it is enough to show that the complementary
probabilities tend to \(0\).  Let \(\mathcal D_n\) be the fiber-degree event
from \(\noderef{FiberAndDegreeEvent}\).  By
\(\noderef{ParameterHierarchy}\), \(0<\varepsilon_2\), and
\(\noderef{FiberAndDegreeConcentration}\), applied with the rounded law
\(\beta=\frac12\) and
\(M(n)=\noderef{TwoBiteNaturalReducedVertexCount}(n)\), the probabilities
\(\Pr(\mathcal D_n)\) tend to \(1\).  Equivalently,
\(\Pr(\mathcal D_n^c)\to0\).  Thus, given \(\rho>0\), there is a threshold
\(N_{\mathcal D}\) such that
\[
\Pr(\mathcal D_n^c)<\rho/2
\]
for all \(n\ge N_{\mathcal D}\).  We shall prove
\[
\Pr(\mathcal M_n^c\cap\mathcal D_n)\to0. \tag{3}
\]
Then there is \(N_{\mathcal B}\) such that
\(\Pr(\mathcal M_n^c\cap\mathcal D_n)<\rho/2\) for all
\(n\ge N_{\mathcal B}\).  The set-theoretic inclusion
\[
\mathcal M_n^c\subseteq
\mathcal D_n^c\cup(\mathcal M_n^c\cap\mathcal D_n)
\]
then gives, by finite subadditivity,
\[
\Pr(\mathcal M_n^c)
\le \Pr(\mathcal D_n^c)+\Pr(\mathcal M_n^c\cap\mathcal D_n)<\rho
\]
for all \(n\ge\max\{N_{\mathcal D},N_{\mathcal B}\}\).  This is the required
\(\noderef{WithHighProbability}\) reduction.

Fix \(n>n_0\) and suppose that \(\mathcal M_n\) fails while
\(\mathcal D_n\) holds.  Negating the universal quantifier in
\(\mathcal M_n\), there is a finite set \(I\subseteq\mathrm{Fin}(n)\) with
\(|I|=K\) for which the left side of (1) is larger than
\(\varepsilon_1K^2\).  Set
\[
Y_I=\sum_{x\in M_I}|X_x(I)|.
\]
By (2),
\[
\left|\bigcup_{x\in M_I}\binom{X_x(I)}2\right|
\le \sum_{x\in M_I}\binom{|X_x(I)|}{2}
\le t_2\,Y_I.
\]
The first inequality is the finite union bound over the witnesses \(x\in M_I\).
For the second, each summand has \(0\le |X_x(I)|\le t_2\), so
\(\binom{|X_x(I)|}{2}\le |X_x(I)|^2/2\le t_2|X_x(I)|\); summing gives the
displayed bound.
Therefore
\[
Y_I>\frac{\varepsilon_1K^2}{t_2}.
\]
The \(n>n_0\) instance of
\(\noderef{ParameterHierarchyEventualComparisons}\) contains
\[
t_2\,n^{1/4-\varepsilon}\le \varepsilon_1K,
\]
and \(t_2=n^{1/4+\varepsilon}>0\), so division by \(t_2\) gives
\[
Y_I>K\,n^{1/4-\varepsilon}. \tag{4}
\]

Choose an inclusion-minimal subfamily \(B\subseteq M_I\) with
\[
\sum_{x\in B}|X_x(I)|>K\,n^{1/4-\varepsilon}.
\]
Such a \(B\) exists because \(M_I\) is finite and the full sum is larger than
the threshold by (4): starting from \(M_I\), repeatedly delete any element
whose deletion preserves the displayed strict inequality; finiteness stops the
process and gives an inclusion-minimal family.  The empty family has sum
\(0\), so this minimal family is nonempty.  If \(b\in B\), inclusion-minimality
gives
\(\sum_{x\in B\setminus\{b\}}|X_x(I)|\le K n^{1/4-\varepsilon}\); adding back
the single term for \(b\) and using the upper bound \(|X_b(I)|\le t_2\) gives
\[
K\,n^{1/4-\varepsilon}
<Y_B:=\sum_{x\in B}|X_x(I)|
\le K\,n^{1/4-\varepsilon}+t_2
\le \frac32K\,n^{1/4-\varepsilon} \tag{5}
\]
for all sufficiently large \(n\).  The last inequality follows from the
rounding lower bound
\[
K\ge (1+\varepsilon)\sqrt{n\log n}
\]
in \(\noderef{ParameterHierarchyEventualComparisons}\): indeed
\[
\frac{t_2}{K n^{1/4-\varepsilon}}
\le
\frac{n^{-1/2+2\varepsilon}}{(1+\varepsilon)\sqrt{\log n}},
\]
and the right side tends to \(0\) because
\(\varepsilon<1/16\).  After increasing the threshold it is at most \(1/2\).
Since every \(x\in B\) is medium and hence
\(|X_x(I)|>t_3=n^{2\varepsilon}\), (5) gives
\[
|B|\le \frac32K\,n^{1/4-3\varepsilon}. \tag{6}
\]

The next step is the multiplicity point.  We do not treat \(Y_B\) itself as a
sum of independent Bernoulli variables.  Instead we pass from the repeated
incidences counted by \(Y_B\) to distinct base-graph edges.  Split
\[
B=B_R\cup B_B,\qquad B_R=B\cap V_R,\qquad B_B=B\cap V_B.
\]
For \(x\in B_R\) and \(u\in X_x(I)\), the definition of
\(\noderef{TwoBiteLiftedNeighborhood}\) says that the red base edge
\(\{x,\pi_R(u)\}\) is present in \(G_R\).  Similarly, for \(x\in B_B\) and
\(u\in X_x(I)\), the blue base edge \(\{x,\pi_B(u)\}\) is present in \(G_B\).
Let \(E_{I,B}\) be the number of distinct present red edges between
\(B_R\) and \(\pi_R(I)\), plus the number of distinct present blue edges
between \(B_B\) and \(\pi_B(I)\).

On \(\mathcal D_n\), every red or blue fiber has size at most
\((1+\varepsilon_2)\log^2 n\): this is the upper side of the
\(\noderef{WithinRelativeError}\) fiber condition in
\(\noderef{FiberAndDegreeEvent}\).  Since
\(\noderef{ParameterHierarchy}\) gives
\(0\le\varepsilon_2<\varepsilon<1/16\), each fiber has size
\(<(17/16)\log^2 n\).

Define a charge from repeated incidences to distinct base edges.  A red
incidence \((x,u)\), with \(x\in B_R\) and \(u\in X_x(I)\), is charged to the
unordered red edge \(\{x,\pi_R(u)\}\).  This edge is present by
\(\noderef{TwoBiteLiftedNeighborhood}\); moreover \(x\ne\pi_R(u)\), since a
loop is never adjacent in the simple red base graph.  Conversely, fix one
present red base edge \(\{a,b\}\).  It can receive charges only from the center
\(a\), through vertices \(u\in I\) with \(\pi_R(u)=b\), if \(a\in B_R\), and
from the center \(b\), through vertices \(u\in I\) with \(\pi_R(u)=a\), if
\(b\in B_R\).  There are no other red centers whose charged edge is
\(\{a,b\}\).  Hence this edge receives at most
\[
|I\cap F_R(a)|+|I\cap F_R(b)|
\le 2(1+\varepsilon_2)(\log n)^2
<3(\log n)^2
\]
red charges.  If only one endpoint lies in \(B_R\), the corresponding single
fiber term is the only possible contribution, so the same upper bound still
holds.  The identical charge argument in the blue graph, using blue fibers and
\(\pi_B\), gives the same bound for each distinct blue base edge.
Therefore
\[
Y_B\le 3(\log n)^2 E_{I,B}. \tag{7}
\]
Combining this with the lower bound in (5) gives
\[
E_{I,B}>
A:=\frac{K n^{1/4-\varepsilon}}{3(\log n)^2}. \tag{8}
\]

Now fix a candidate pair \((I,B)\) satisfying \(|I|=K\) and
\(|B|\le\frac32K n^{1/4-3\varepsilon}\), and condition only on the injection
coordinate.  If no injection is compatible with the sample parameters, then
the conditional fiber is empty and the fixed-pair probability is \(0\), so
assume an injection has been fixed.  By \(\noderef{TwoBiteSample}\), this
fixes the fibers, the projections \(\pi_R(I)\) and \(\pi_B(I)\), and the
red-blue split of \(B\).  The remaining randomness is exactly the pair of base
graphs.  The factorization in \(\noderef{TwoBiteSampleWeight}\) as
\(\noderef{UniformInjectionWeight}\) times the two
\(\noderef{GnpGraphWeight}\) factors says that, after the injection is fixed,
the red base-edge indicators are mutually independent Bernoulli variables with
parameter \(p\), the blue base-edge indicators are mutually independent
Bernoulli variables with the same parameter, and the red and blue families are
independent of each other.  Thus the conditional law is the product law used in
the binomial domination below, with no further conditioning on
\(\mathcal D_n\).

For this fixed injection, let \(T_R(I,B)\) be the set of unordered nonloop red
pairs \(\{x,y\}\) with \(x\in B_R\) and \(y\in\pi_R(I)\), and let
\(T_B(I,B)\) be the analogous blue set with \(B_B\) and \(\pi_B(I)\).  Then
\[
|T_R(I,B)|+|T_B(I,B)|
\le |B_R|\,|\pi_R(I)|+|B_B|\,|\pi_B(I)|
\le |B|K.
\]
The count \(E_{I,B}\) is at most the sum of the independent Bernoulli edge
indicators over \(T_R(I,B)\cup T_B(I,B)\).  Loops have been omitted, and if two
descriptions name the same unordered edge they contribute only one element to
\(T_R(I,B)\) or \(T_B(I,B)\); both operations can only decrease this sum.
Adding independent dummy Bernoulli-\(p\) variables until there are exactly
\(|B|K\) trials gives a binomial random variable that dominates \(E_{I,B}\)
for this injection.  Consequently \(E_{I,B}\) is stochastically dominated,
conditionally on the injection, by a binomial variable
\[
Z_{I,B}\sim\mathrm{Bin}(|B|K,p),
\]
where
\[
p=\noderef{TwoBiteEdgeProbability}\!\left(\frac12,n\right)
=\frac12\sqrt{\frac{\log n}{n}}.
\]
Equivalently, if \(\mathcal B_{I,B}\) denotes the event
\(\mathcal D_n\cap\{E_{I,B}>A\}\) for this fixed pair, then after conditioning
on any injection \(\pi\) we may bound \(\Pr(\mathcal B_{I,B}\mid\pi)\) by the
upper tail of \(Z_{I,B}\); the condition \(\mathcal D_n\) is not needed for
this probability upper bound, because intersecting with it only removes
samples.
Its mean \(\mu_{I,B}\) satisfies
\[
\mu_{I,B}\le p|B|K
\le \frac32pK^2 n^{1/4-3\varepsilon}
\le 2pK^2 n^{1/4-3\varepsilon}. \tag{9}
\]
The rounded-scale part of
\(\noderef{ParameterHierarchyEventualComparisons}\) gives
\[
K<(1+\varepsilon)\sqrt{n\log n}+1.
\]
After increasing the threshold, the summand \(1\) is at most
\((1+\varepsilon)\sqrt{n\log n}\), and therefore
\[
K\le 2(1+\varepsilon)\sqrt{n\log n}.
\]
Combining this with (9) and (8), we obtain
\[
\frac{\mu_{I,B}}{A}
\le 6pK(\log n)^2 n^{-2\varepsilon}
\le 6(1+\varepsilon)(\log n)^3 n^{-2\varepsilon}. \tag{10}
\]
Since \(\varepsilon>0\), the right side tends to \(0\).  Increase the threshold
so that it is at most \(1/(4e)\).  Then
\(\mu_{I,B}\le A/(4e)\), in particular \(\mu_{I,B}\le A/2\).  Applying a finite union bound gives the upper tail. The event \(Z_{I,B}\ge A\) requires at least \(t=\lceil A\rceil\) of the \(|B|K\) independent Bernoulli-\(p\) candidate edges to be present. There are at most \(\binom{|B|K}{t}\) subsets of \(t\) candidates, and for each subset the probability that all its edges are present is \(p^t\). Hence
\[
\Pr(Z_{I,B}\ge A)
\le \binom{|B|K}{t}p^t
\le \left(\frac{e |B| K p}{t}\right)^t
= \left(\frac{e\mu_{I,B}}{t}\right)^t
\le \left(\frac{e\mu_{I,B}}{A}\right)^A .
\]
The preceding inequality \(\mu_{I,B}\le A/(4e)\) makes the base at most
\(1/4\), so this is at most \(\exp(-A\log 4)\), and hence at most
\(\exp(-A/4)\).  Therefore
\[
\Pr(E_{I,B}>A\mid\pi)
\le \Pr(Z_{I,B}>A)
\le \exp(-A/4). \tag{11}
\]
The lower rounded-scale bound
\[
K\ge(1+\varepsilon)\sqrt{n\log n}
\]
also gives
\[
A
\ge \frac{1+\varepsilon}{3}\,
  \frac{n^{3/4-\varepsilon}}{(\log n)^{3/2}}.
\]
Increase the threshold again so that
\[
\frac{1+\varepsilon}{3}\,
  \frac{n^\varepsilon}{(\log n)^{3/2}}\ge 4.
\]
Then
\[
A=\frac{K n^{1/4-\varepsilon}}{3(\log n)^2}
\ge 4 n^{3/4-2\varepsilon}
\]
for all remaining \(n\).  Hence every fixed candidate \((I,B)\), after
averaging over the injection coordinate as well, satisfies
\[
\Pr(\mathcal B_{I,B})
\le \exp(-n^{3/4-2\varepsilon}). \tag{12}
\]
Indeed, \(\noderef{TwoBiteEventProbability}\) decomposes as the finite
\(\noderef{UniformInjectionWeight}\)-weighted average of the conditional
probabilities over injections.  The estimate (11) is uniform for every fixed
injection, and \(\mathcal D_n\) is only being intersected with the bad event,
so dropping \(\mathcal D_n\) can only increase each conditional probability.

It remains to count choices.  There are at most
\[
\binom{n}{K}\le n^K=\exp(K\log n)
\]
choices for \(I\).  By (6), put
\[
S=K n^{1/4-3\varepsilon},
\qquad
R=\frac32S;
\]
then the witness family \(B\) has at most \(R\) vertices.  The base vertex
universe has \(2M\) vertices, and the rounding comparison in
\(\noderef{ParameterHierarchyEventualComparisons}\) gives
\[
M\le \frac{n}{(\log n)^2}.
\]
After increasing the threshold so that \((\log n)^2\ge2\), this universe has
at most \(n\) vertices.  Let \(r=\lceil R\rceil\).  Every witness family
satisfying (6) has some cardinality \(j\le r\), and for each such \(j\) there
are at most \(n^j\) choices.  Hence the number of possible \(B\)'s is at most
\[
\sum_{0\le j\le r} n^j\le (r+1)n^r.
\]
Increase the threshold so that \(S\ge2\).  Since \(R=\frac32S\), this gives
\(r\le R+1=\frac32S+1\).  Also
\[
\frac{\log(2S+1)}{S\log n}\to0,
\]
because \(S\) is at least a positive constant times
\(n^{3/4-3\varepsilon}\sqrt{\log n}\) and \(\varepsilon<1/16\).  Increase the
threshold again so that both \(S\ge4\) and
\[
\log n+\log(2S+1)\le\frac12S\log n.
\]
This is possible because \(1/S\to0\) and
\(\log(2S+1)/(S\log n)\to0\).  Since \(r+1\le2S+1\), we obtain
\[
(r+1)n^r
\le
\exp\!\left(\log(2S+1)+\left(\frac32S+1\right)\log n\right)
\le \exp\!\left(2K n^{1/4-3\varepsilon}\log n\right)
\]
possible choices for \(B\), since \(S=K n^{1/4-3\varepsilon}\).  This bound
counts all color choices and all subsets of the whole red-blue base vertex
universe, so it also covers the particular split \(B=B_R\cup B_B\).
The deterministic construction above gives the required inclusion for the
union bound.  If \(\mathcal M_n^c\cap\mathcal D_n\) occurs, choose an
offending \(K\)-set \(I\) and then choose \(B\subseteq M_I\) minimal with
\(\sum_{x\in B}|X_x(I)|>K n^{1/4-\varepsilon}\).  Equations (6) and (8) show
that this pair is among the counted candidates and satisfies
\(E_{I,B}>A\).  Hence
\[
\mathcal M_n^c\cap\mathcal D_n
\subseteq
\bigcup_{\substack{|I|=K\\ |B|\le R}}
\mathcal B_{I,B}.
\]
No disjointness is asserted or needed.  Combining this inclusion, the count of
candidates, and the fixed-pair probability bound (12), the union bound gives
\[
\Pr(\mathcal M_n^c\cap\mathcal D_n)
\le
\exp\!\left(
K\log n+2K n^{1/4-3\varepsilon}\log n
-n^{3/4-2\varepsilon}
\right). \tag{13}
\]
This is exactly the medium exponential comparison included in the
\(n>n_0\) instance of
\(\noderef{ParameterHierarchyEventualComparisons}\):
\[
\exp\!\left(
K\log n+2K n^{1/4-3\varepsilon}\log n
-n^{3/4-2\varepsilon}
\right)\le\varepsilon_1.
\]
Thus all rounded parameters \(M,K,p,t_3\) and the finite hierarchy comparison
used in the medium union bound are the ones supplied by the child hierarchy
node.

For the limit required in (3), we also record the explicit decay of the same
left-hand side.  The rounded upper bound used above gives, for all sufficiently
large \(n\),
\[
K\le 2(1+\varepsilon)\sqrt{n\log n}.
\]
Consequently
\[
\frac{K\log n}{n^{3/4-2\varepsilon}}
\le
2(1+\varepsilon)(\log n)^{3/2}n^{-1/4+2\varepsilon}
\]
and
\[
\frac{2K n^{1/4-3\varepsilon}\log n}{n^{3/4-2\varepsilon}}
\le
4(1+\varepsilon)(\log n)^{3/2}n^{-\varepsilon}.
\]
Both right sides tend to \(0\): the first because
\(-1/4+2\varepsilon<0\), and the second because \(\varepsilon>0\).  Increase
the threshold so that their sum is at most \(1/2\).  Then the two positive
terms in the exponent of (13) contribute at most
\(\frac12 n^{3/4-2\varepsilon}\), so that exponent is at most
\[
-\frac12 n^{3/4-2\varepsilon},
\]
and hence
\[
\Pr(\mathcal M_n^c\cap\mathcal D_n)
\le \exp\!\left(-\frac12 n^{3/4-2\varepsilon}\right)\to0.
\]
The last convergence uses \(\varepsilon<1/16\), hence
\(3/4-2\varepsilon>0\), so \(n^{3/4-2\varepsilon}\to\infty\).
This proves (3).  By the initial \(\noderef{WithHighProbability}\) reduction,
the probability of \(\mathcal M_n\) tends to \(1\).  This is exactly the
claimed \(\noderef{WithHighProbability}\) statement under the rounded
\(M(n)\), edge probability \(p(n)\), and scale \(K(n)\).
\end{proof}
